\section{Specialization and a missing spectral component}
\label{sec:specialization}

We first translate a missing component into a geometric vanishing statement.
The object we obtain will be compact, hence constructible on finite-type
charts.  More importantly, its Hecke transforms will have zero nearby cycles
after any pullback whose projection to the bundle stack is smooth.  The map
recording the Hecke point need not be smooth.  This last qualification is
what permits the hypersurface tests in Sections~\ref{sec:twohecke}
and~\ref{sec:slices}.

Fix data satisfying Hypothesis~\ref{hyp:A} or Hypothesis~\ref{hyp:B}.
Nothing in this section assumes the full-support conclusion of
Theorem~\ref{thm:full-support}.

\subsection{A trait and geometric specialization}

Initially set
\[
 R=W(k),\qquad S=\Spec(R),\qquad
 K=\overline{\operatorname{Frac}(R)}.
\]
We fix the algebraic closure $K$ throughout.  Since $k$ is perfect, every
nonzero Witt vector is a power of $p$ times a unit.  Thus $R$ is a complete
Noetherian discrete valuation ring with uniformizer $p$ and residue field
$k$.  It is strictly henselian, and it is excellent because it is a complete
Noetherian local ring; see \cite[Tag~07QS]{Stacks}.

Choose a smooth projective lift $\mathscr X\to S$ of $X$.
For completeness, such a lift exists over $W(k)$ for every algebraically
closed $k$ in either hypothesis.  At each infinitesimal lifting step the
obstruction for a smooth curve lies in $H^2(X,T_X)$, which vanishes.
An ample invertible sheaf also lifts, since its obstruction lies in
$H^2(X,\mathcal O_X)=0$.  Algebraization of the resulting proper formal
curve with an ample invertible sheaf gives a projective lift.  It is smooth
near the special fiber; the nonsmooth locus, being closed in a proper scheme
over the local base, would meet the special fiber if it were nonempty.
Hence the lift is smooth everywhere.  These deformation and algebraization
statements can be found in \cite[Tags~0DZQ and~0E7R]{Stacks}; the same lift is
used in \cite[\S3.1.1]{GR}.  This construction requires no finite-field
descent of $X$.

The root datum of the split group $G$ gives a split reductive group over
$\mathbb Z$ and hence over $S$
\cite[Theorems~6.1.16(2) and~6.1.17]{Conrad}; we retain the notation
$G$ for this lift.
Write
\[
 \mathcal B=\Bun_G(\mathscr X/S),\qquad
 s=\Spec(k),\qquad \eta=\Spec(K).
\]
The stack $\mathcal B=\Hom_S(\mathscr X,B_SG)$ is algebraic and locally
of finite presentation over $S$, with affine diagonal, by
\cite[Theorem~1.2(i)--(iii)]{HallRydh}. Indeed, $\mathscr X/S$ is proper,
flat and finitely presented, while the classifying stack $B_SG$ is
locally of finite presentation and quasi-separated, with affine
stabilizers and affine diagonal. The obstruction group for deforming
a bundle on a curve is the second cohomology of its adjoint bundle,
which vanishes. Thus $\mathcal B$ is smooth over $S$.

We will replace $R$ by its integral closure in a finite extension of
$\operatorname{Frac}(R)$ inside $K$ when necessary.  Such an integral closure
is a finite complete discrete valuation ring with residue field $k$;
see \cite[Tags~0EXQ and~0323]{Stacks}.  It remains excellent and strictly
henselian.  We retain $R,S,\mathscr X$ and $\mathcal B$ for the resulting
models.  All dimension arguments below take place over these Noetherian
traits, not over the integral closure in the entire field $K$.

For a finite-type scheme or algebraic space $T$ over $S$, let
\[
 \Psi_T:\Shv(T_K)\longrightarrow\Shv(T_s)
\]
be geometric specialization, with monodromy forgotten.  Here, on a
finite-type scheme over a field, $\Shv$ is the ind-completion of the
category of bounded constructible $E$-adic complexes.  We use ordinary
$*$-pullbacks, ordinary tensor products and ordinary stalks.  The same
notation applies to locally finite-type algebraic stacks by smooth descent.
In particular, the functor $\Psi_{\mathcal B}$ is defined on the entire
sheaf category.

We recall how the construction in \cite[\S\S4.1.1--4.1.5]{GR} fixes the
meaning of this functor.  On a finite-type affine chart one first uses
geometric nearby cycles for finite coefficient rings at finite extensions
of the generic field.  One then passes through the compatible coefficient
limits, inverts $\ell$, extends the coefficient field to $E$, and extends
continuously from constructible objects.  Compatibility with smooth
pullback gives the stack version.  In particular, $\Psi_T$ is an exact
functor of stable categories, is continuous, and preserves constructibility.
It commutes with smooth pullback and representable proper pushforward
\cite[\S4.1.6]{GR}.

These functors are unchanged by the allowed finite trait replacements,
using the identified geometric fibers.  Indeed, after fixing one finite
extension, the finite subextensions of $K$ containing it form a cofinal
subsystem: enlarge any stage by taking its compositum with the fixed
extension.  The models and torsion nearby-cycle functors at the remaining
stages are the same.  All subsequent coefficient operations and
ind-extensions therefore agree.  For $T=S$, the functor $\Psi_S$ is the
identity on the underlying complexes of geometric-point coefficients.
At no stage do we take inertia invariants.

We will also use the following tensor compatibility.  If
$a:T\to\mathcal B$ is smooth and $\mathcal L$ is a finite-rank lisse sheaf
on $T$, then for any $D\in\Shv(\mathcal B_K)$,
\begin{equation}
 \Psi_T\bigl(a_K^*D\otimes\mathcal L_K\bigr)
 \simeq
 a_s^*\Psi_{\mathcal B}(D)\otimes\mathcal L_s.
 \label{eq:spec-lisse-tensor}
\end{equation}
A lisse sheaf on the source of a smooth morphism is universally locally
acyclic over the target.  Thus \eqref{eq:spec-lisse-tensor} is the
universally locally acyclic tensor-pullback formula of
\cite[Lemma~4.1.9]{GR}.  Equivalently, it follows from smooth base change
and the tensor formula for a lisse sheaf extending over the model.
Continuity makes the formula applicable to arbitrary $D$ in the indicated
ind-category.

\subsection{The spectral summand and compactness}

Let
\[
 \mathcal C_s=\Shv_{\Nilp}(\mathcal B_s),\qquad
 \mathcal C_K=\Shv_{\Nilp}(\mathcal B_K),\qquad
 Y_K=LS^{\mathrm{restr}}_{\check G}(\mathscr X_K).
\]
Recall that $Y=LS^{\mathrm{restr}}_{\check G}(X)$ denotes the special-fiber
spectral stack. Restriction to the special fiber gives an equivalence
$\QLisse(\mathscr X)\simeq\QLisse(X)$
\cite[\S3.1.3]{GR}. Composing its inverse with restriction to the
geometric generic fiber defines the symmetric monoidal functor
\[
 \rho:\QLisse(X)\longrightarrow\QLisse(\mathscr X_K).
\]
The induced map
\[
 i:Y\hookrightarrow Y_K
\]
identifies $Y$ with an open-and-closed union of connected components
\cite[\S3.1.3]{GR}.  In particular,
\begin{equation}
 \mathcal D_Y
 :=\QCoh(Y)\underset{\QCoh(Y_K)}\otimes\mathcal C_K
 \label{eq:spec-DY}
\end{equation}
is a direct summand of $\mathcal C_K$.

We need the compatibility of \emph{actual} specialization with this
summand.  Corollary~4.3.6 and \S4.3.9 of \cite{GR} show that the restriction
of the chartwise functor $\Psi_{\mathcal B}$ constructed above gives
\[
 \Psi_{\mathcal B}:\mathcal D_Y\longrightarrow\mathcal C_s.
\]
It is $\QCoh(Y)$-linear by \cite[Theorem~3.1.6(B) and \S4.3.10]{GR}.
Thus the functor used here is not merely one obtained by transporting a
spectral functor across Langlands equivalences.  Its values on smooth charts
are the geometric nearby cycles used later.

The characteristic-zero equivalence of \cite[Theorem~1.3.9(ii)]{GR} and
its spectral linearity identify
\begin{equation}
 \mathcal D_Y\simeq\IndCoh_{\Nilp}(Y).
 \label{eq:spec-generic-equivalence}
\end{equation}
In positive characteristic we use only the given equivalence onto the
primed union
\[
 \mathcal C_s\simeq\IndCoh_{\Nilp}(Y'),\qquad Y'\subseteq Y.
\]

\begin{proposition}
\label{spec:missing}
Suppose a connected component $C\subset Y$ is disjoint from $Y'$.
Then
\[
 \mathcal D_C
 :=\QCoh(C)\underset{\QCoh(Y)}\otimes\mathcal D_Y
\]
is nonzero, and $\Psi_{\mathcal B}(D)=0$ for every
$D\in\mathcal D_C$.  Moreover, there exists a nonzero
$F\in\mathcal D_C$ that is compact both in $\mathcal C_K$ and in
$\Shv(\mathcal B_K)$.  For every finite-type smooth scheme chart
$b:U\to\mathcal B$, the complex $b_K^*F$ is bounded constructible.
\end{proposition}

\begin{proof}
Let $e_C\in\QCoh(Y)$ be the structure sheaf of $C$ extended by zero.
Its action is the projection onto $\mathcal D_C$.  On $\mathcal C_s$ this
object acts by zero, since $C\cap Y'=\varnothing$ and the primed equivalence
is spectral-linear.  Hence, for $D\in\mathcal D_C$,
\[
 \Psi_{\mathcal B}(D)
 \simeq\Psi_{\mathcal B}(e_C\cdot D)
 \simeq e_C\cdot\Psi_{\mathcal B}(D)=0.
\]
By \eqref{eq:spec-generic-equivalence}, the category $\mathcal D_C$ is
identified with $\IndCoh_{\Nilp}(C)$.  It is nonzero: the fully faithful
functor from $\QCoh(C)$ into $\IndCoh_{\Nilp}(C)$ sends the nonzero object
$\mathcal O_C$ to a nonzero object; see \cite[\S\S1.3.1--1.3.4]{GR}.

The category $\mathcal D_Y$ is compactly generated
\cite[\S3.1.7]{GR}.  The binary decomposition
\[
 \mathcal D_Y\simeq\mathcal D_C\times\mathcal D_{Y\setminus C}
\]
shows that projection onto $\mathcal D_C$ preserves compactness: its right
adjoint, the inclusion of that factor, is continuous.  Project compact
generators of $\mathcal D_Y$.  At least one projection is nonzero, because
otherwise their colimit closure would give $\mathcal D_C=0$.  Choose such
an object $F$.

The inclusion of either direct factor also preserves compactness, since
its right adjoint is the continuous projection.  Applying this first to
$\mathcal D_C\subset\mathcal D_Y$ and then to
$\mathcal D_Y\subset\mathcal C_K$ shows that $F$ is compact in
$\mathcal C_K$.  The nilpotent inclusion
$\mathcal C_K\hookrightarrow\Shv(\mathcal B_K)$ preserves compactness by
\cite[Theorem~1.1.7]{GR}.  Finally, a compact sheaf on an algebraic stack
is constructible on smooth finite-type charts
\cite[\S F.2.2]{AGKRRV}.  This proves the last assertion.
\end{proof}

The argument uses one missing component and a binary direct-factor
decomposition.  It does not require that $Y$ have finitely many components
or that the unit $e_C$ be compact.  If $Y'\ne Y$, such a component $C$
exists because $Y'$ is a union of the open-and-closed components of $Y$.
For the rest of the geometric proof, fix $C$ and $F$ as in
Proposition~\ref{spec:missing}.

\subsection{Hecke vanishing with an arbitrary leg}

Let $V$ be a finite-dimensional representation of $\check G$, and let
\[
 \mathrm H_V:\mathcal C_K\longrightarrow
       \Shv(\mathcal B_K\times_K\mathscr X_K)
\]
be the one-leg Hecke functor, with the normalization giving its usual
spectral-action description.  The correspondence and its kernel will be
specified in Section~\ref{sec:hecke}.  Here we use the universal evaluation
objects
\[
 \mathcal E_{V,s}\in\QCoh(Y)\otimes\QLisse(X),\qquad
 \mathcal E_{V,K}\in\QCoh(Y_K)\otimes\QLisse(\mathscr X_K).
\]
For an affine test scheme mapping to a spectral stack, the restriction of
$\mathcal E_V$ is the family of local systems obtained by applying the
corresponding tensor functor to $V$.  Acting by its first tensor factor and
then taking the exterior product gives the Hecke functor
\cite[\S14.3.3]{AGKRRV}.

The definition of $i$ implies
\begin{equation}
 (i^*\otimes\id)(\mathcal E_{V,K})
 \simeq (\id\otimes\rho)(\mathcal E_{V,s}).
 \label{eq:spec-universal-family}
\end{equation}
Indeed, this identity holds on every affine test mapping to $Y$: both sides
are the local system obtained by extending the specified special-fiber
family.  To pass from these tests to the global identity, note that
$\QLisse$ of a smooth curve is compactly generated and hence dualizable
as a presentable dg category
\cite[\S1.2.11 and \S\S E.2.6--E.2.7]{AGKRRV}.
Tensoring with a dualizable category preserves limits, so it commutes with
the descent limit defining $\QCoh(Y)$.  This justifies the passage from the
affine identities to \eqref{eq:spec-universal-family}.

A second, distinct input is
\begin{equation}
 \mathcal E_{V,s}\in\QCoh(Y)\otimes\IndLisse(X),
 \qquad \IndLisse(X)=\operatorname{Ind}(\operatorname{Lisse}(X)),
 \label{eq:spec-IndLisse}
\end{equation}
proved in \cite[Remark~14.3.9]{AGKRRV}.  Here
$\operatorname{Lisse}(X)$ consists of bounded constructible complexes
whose ordinary cohomology sheaves are finite-rank lisse sheaves.
We do not identify $\IndLisse(X)$ with $\QLisse(X)$: that identification
fails for $X=\mathbb P^1$.  The dualizability used in descent and the
stronger membership assertion \eqref{eq:spec-IndLisse} have different roles.

Applying the $\QCoh(Y)$-action in \eqref{eq:spec-IndLisse} to $F$ gives an
object of $\mathcal D_C\otimes\IndLisse(X)$.  After applying
$\id\otimes\rho$ and the exterior-product functor, its image is $\mathrm H_V(F)$
by \eqref{eq:spec-universal-family}.  Elementary tensors generate the tensor
product of presentable dg categories under colimits.  Bounded truncations
also build every object of $\operatorname{Lisse}(X)$ from finite-rank
lisse sheaves.  Consequently $\mathrm H_V(F)$ belongs to the stable colimit closure
of objects
\begin{equation}
 D\boxtimes\mathcal L_K,\qquad D\in\mathcal D_C,
 \label{eq:spec-exterior-generators}
\end{equation}
where $\mathcal L_K$ is the restriction of a finite-rank lisse sheaf
$\mathcal L_S$ on $\mathscr X$.
To see the extension assertion, proper henselian invariance identifies the
finite \emph{\'etale} covers of $\mathscr X$ with those of $X$
\cite[Tag~0A48]{Stacks}.  Apply this to the compatible finite coefficient
local systems of a lisse lattice, then invert $\ell$ and extend coefficients.
The resulting extension is the one inducing $\rho$.

\begin{lemma}
\label{spec:vanishing}
Let $a:T\to\mathcal B$ be smooth, where $T$ is a scheme or algebraic space
of finite type over $S$, and let $\ell_T:T\to\mathscr X$ be any $S$-morphism.
For $F$ as in Proposition~\ref{spec:missing} and every finite-dimensional
representation $V$ of $\check G$, one has
\[
 \Psi_T(a_K^*F)=0,
 \qquad
 \Psi_T\bigl((a,\ell_T)_K^*\mathrm H_V(F)\bigr)=0.
\]
The same statements hold on the smooth locus of $a$, and after any of the
finite trait replacements described above.
\end{lemma}

\begin{proof}
The first equality follows from smooth base change and
$\Psi_{\mathcal B}(F)=0$.
For an exterior product in \eqref{eq:spec-exterior-generators}, ordinary
pullback gives
\[
 (a,\ell_T)_K^*(D\boxtimes\mathcal L_K)
 =a_K^*D\otimes\ell_{T,K}^*\mathcal L_K.
\]
The sheaf $\ell_T^*\mathcal L_S$ is lisse on $T$, regardless of the
smoothness of $\ell_T$.  Formula~\eqref{eq:spec-lisse-tensor} and
Proposition~\ref{spec:missing} therefore give
\[
 \Psi_T\bigl(a_K^*D\otimes\ell_{T,K}^*\mathcal L_K\bigr)
 \simeq a_s^*\Psi_{\mathcal B}(D)\otimes\ell_{T,s}^*\mathcal L_s=0.
\]
Both pullback and specialization are continuous, so the same vanishing
holds for the stable colimit closure of these exterior products, which
contains $\mathrm H_V(F)$.  The calculation is local on the smooth locus, and the
finite-extension assertion follows from the invariance of geometric
specialization already proved.
\end{proof}

We have now obtained the two vanishings needed for the geometric argument.
Their proof uses smoothness only of the map to the bundle stack; it makes
no smooth-base-change assertion for the combined map $(a,\ell_T)$.
The next two sections construct modifications that turn a non-nilpotent
covector into a test to which Lemma~\ref{spec:vanishing} applies.
